\selectlanguage{english}

\section{Metric manifolds}

\begin{definition}{Metric}{metric}
  A \bcdef{metric} $ g $ is a (0,2) tensor field on a manifold $ \mathcal{M} $ that is:
  \begin{enumerate}
    \item symmetric: $ g(X,Y) = g(Y,X) $;
    \item non-degenerate: $ \exists p \in \mathcal{M} : g(X,Y)\vert_p = 0 \,\,\forall Y \in T_p \mathcal{M} \iff X_p = 0 $.
  \end{enumerate}
  A manifold equipped with a metric $ (\mathcal{M} , g) $ is called a \bcdef{metric manifold}.
\end{definition}

With a choice of coordinates, the metric can be written as:
\begin{equation}
  g = g_{\mu \nu}(x) \dd x^{\mu} \otimes \dd x^{\nu}
\end{equation}
where:
\begin{equation}
  g_{\mu \nu} = g\left( \frac{\pa}{\pa x^{\mu}}, \frac{\pa}{\pa x^{\nu}} \right)
\end{equation}
It is often written also as $ \dd s^2 = g_{\mu \nu}(x) \dd x^{\mu} \dd x^{\nu} $. The matrix $ g_{\mu \nu}(x) \in \R^{n \times n} $ is symmetric, and there's always a choice of basis on each tangent space such that this matrix is diagonal: the non-degeneracy condition implies that none of the diagonal elements vanish.

\begin{proposition}{Signature}{}
  The \bcprop{signature} of a metric, i.e. the number of negative entries when diagonalized, is independent on the choice of basis.
\end{proposition}

\begin{proofbox}
  \begin{proof}
    From Sylvester's theorem of inertia.
  \end{proof}
\end{proofbox}

\subsection{Riemannian manifolds}

\begin{definition}{Riemannian manifold}{}
  A \bcdef{Riemannian manifold} $ (\mathcal{M},g) $ is a manifold equipped with a metric with totally-positive signature.
\end{definition}

\begin{example}{Euclidean space}{}
  The Euclidean space $ \R^n $, equipped with the metric $ g_{\mu \nu} = \delta_{\mu \nu} $ (in Cartesian coordinates), is a Riemannian manifold.
\end{example}

In a Riemannian manifold $ (\mathcal{M},g) $ and $ X \in \xm $, the \bctxt{length} of $ X $ at $ p \in \mathcal{M} $ is defined as:
\begin{equation}
  \abs{X_p} \defeq \sqrt{g(X,X)\vert_p}
\end{equation}
Given $ Y \in \xm $, the \bctxt{angle} between $ X $ and $ Y $ at $ p \in \mathcal{M} $ is defined as:
\begin{equation}
  \cos \theta \defeq \frac{g(X,Y) \vert_p}{\abs{X_p} \abs{Y_p}}
\end{equation}
This can be generalized to distances between points on a curve $ \sigma : \mathcal{\R} \rightarrow \mathcal{M} $:
\begin{equation}
  d(p,q) = \int_a^b \dd t \sqrt{g(X,X)\vert_{\sigma(t)}}
  \label{eq:distance}
\end{equation}
where $ \sigma(a) = p $, $ \sigma(b) = q $ and $ X $ is the tangent vector field of the curve. With parametrization $ x^{\mu}(t) $, the tangent vector has components $ X^{\mu} = \frac{\dd x^{\mu}}{\dd t} $, thus:
\begin{equation}
  d(p,q) = \int_a^b \dd t \sqrt{g_{\mu \nu}(x) \frac{\dd x^{\mu}}{\dd t} \frac{\dd x^{\nu}}{\dd t}}
\end{equation}
It is important to note that this distance is independent of the parametrization.

\subsection{Lorentzian manifolds}

\begin{definition}{Lorentzian manifold}{}
  A \bcdef{Lorentzian manifold} $ (\mathcal{M},g) $ is a manifold equipped with a metric which has a signature with a single negative sign.
\end{definition}

\begin{example}{Lorentz--Minkowski metric}{}
  The simplest Lorentzian manifold is $ \R^n $ with the \bcex{Lorentz--Minkowski metric}:
  \begin{equation}
    \eta = - \dd x^0 \otimes \dd x^0 + \dd x^1 \otimes \dd x^1 + \dots + \dd x^{n-1} \otimes \dd x^{n-1}
  \end{equation}
  Its components are $ \eta_{\mu \nu} = \diag (-1,+1,\dots,+1) $, thus this is a Lorentzian manifold.
\end{example}

On a general Lorentzian manifold, at any point $ p \in \mathcal{M} $ it is always possible to choose an orthonormal basis $ \{e_{\mu}\}_{\mu = 0,\dots,n-1} $ of $ T_p \mathcal{M} $ such that $ g_{\mu \nu} \vert_p = \eta_{\mu \nu} $: this fact is closely related to the equivalence principle. Consider a different basis $ \tilde{e}_{\mu} = \tensor{\Lambda}{^\nu _\mu} e_{\nu} $: the condition for it to leave the Lorentz--Minkowski metric unchanged is:
\begin{equation}
  \eta_{\mu \nu} = \tensor{\Lambda}{^\rho _\mu} \tensor{\Lambda}{^\sigma _\nu} \eta_{\rho \sigma}
\end{equation}
This is the defining equation of a Lorentz transformation: on a Lorentzian manifold, the basic features of special relativity are locally recovered. Thus, other ideas from special relativity can be imported.

\begin{definition}{Curves in a Lorentzian manifold}{}
  Given a Lorentzian manifold $ (\mathcal{M},g) $ and $ X \in \xm $, at $ p \in \mathcal{M} $ the vector field is said to be \bcdef{timelike} if $ g(X_p,X_p) < 0 $, \bcdef{null} if $ g(X_p,X_p) = 0 $ or \bcdef{spacelike} if $ g(X_p,X_p) > 0 $.
\end{definition}

At each point $ p \in \mathcal{M} $ it is possible to draw \bctxt{lightcones}, i.e. the null tangent vectors at that point, which are past-directed or future-directed: these lightcones vary smoothly as the point is varied smoothly on the manifold, elucidating the causal structure of spacetime.

The distance between two points on a curve depends on the nature of the tangent vector field of the curve: a \textit{timelike curve} is a curve whose tangent vector field is everywhere timelike, and analogously for the other cases. The distance on a spacelike curve is defined as in \eref{eq:distance}, while that on a timelike curve gets a negative sign in the square root. With parametrization $ x^{\mu}(t) $, it is possible to define the \bctxt{proper time} on a timelike curve as:
\begin{equation}
  \tau = \int_a^b \dd t \sqrt{- g_{\mu \nu}(x) \frac{\dd x^{\mu}}{\dd t} \frac{\dd x^{\nu}}{\dd t}}
\end{equation}
This is precisely the action of a free particle moving in spacetime.

\newpage
\section{Connections}

There is a different way to differentiate tensor fields, distinct from the Lie derivative, associated to a different way to map different vector spaces at different points: the covariant derivative.

For the rest of this chapter, $ \mathcal{M} $ is implied to be an $ n $-dimensional metric manifold with metric tensor $ g $.

\subsection{Covariant derivative}

\begin{definition}{Connection}{}
  The \bcdef{connection} is a map $ \na : \xm \times \xm \rightarrow \xm $, usually written as $ \na(X,Y) \equiv \na_X Y $, where $ \na_X $ is called the \bcdef{covariant derivative}, satisfying the following properties for all $ X,Y,Z \in \xm $ and $ f,g \in \cm $:
  \begin{enumerate}
    \item $ \na_X (Y + Z) = \na_X Y + \na_X Z $;
    \item $ \na_{fX + gY} Z = f \na_X Z + g \na_Y Z $;
    \item $ \na_X (fY) = f \na_X Y + X(f) Y $.
  \end{enumerate}
\end{definition}

Usually $ X(f) \equiv \na_X f $. The covariant derivative endows the manifold with more structure: in particular, given a basis $ \{e_{\mu}\}_{\mu = 1 , \dots , n} $ of $ \xm $, its covariant derivative is expressed as:
\begin{equation}
  \na_{e_{\rho}} e_{\nu} \equiv \Gamma^{\mu}_{\rho \nu} e_{\mu}
\end{equation}
The $ \Gamma^{\mu}_{\rho \nu} $ are the components of the connection on that basis. Conventionally $ \na_{e_{\mu}} \equiv \na_{\mu} $, thus resembling a partial derivative. To elucidate how the covariant derivative acts on vector fields:
\begin{equation*}
  \begin{split}
    \na_X Y
    &= \na_X (Y^\mu e_\mu) \\
    &= X(Y^\mu) e_\mu + Y^\mu \na_X e_\mu \\
    &= X^{\nu} e_{\nu} (Y^{\mu}) e_{\mu} + Y^{\mu} X^{\nu} \na_{\nu} e_{\mu} \\
    &= X^{\nu} \left[ e_{\nu}(Y^{\mu}) + \Gamma^{\mu}_{\nu \rho} Y^{\rho} \right] e_{\mu} \\
    &= X^{\nu} \na_{\nu} Y = X^{\nu} (\na_{\nu} Y)^{\mu} e_{\mu}
  \end{split}
\end{equation*}
The dependency on $ X $ can therefore be eliminated, and in components:
\begin{equation}
  (\na_\nu Y)^\mu = e_\nu (Y^\mu) + \Gamma^\mu_{\nu \rho} Y^\rho
\end{equation}
A sloppy notation is often used: $ (\na_\nu Y)^\mu \equiv \na_\nu Y^\mu $. This must not be confused as the covariant derivative of $ Y^\mu $. Moreover $ \na_\nu Y^\mu \equiv \tensor{Y}{^\mu_{;\nu}} $, while $ \pa_\mu f \equiv f_{,\mu} $. On the coordinate basis $ e_\mu = \pa_\mu $, then:
\begin{equation}
  \tensor{Y}{^\mu_{;\nu}} = \tensor{Y}{^\mu_{,\nu}} + \Gamma^\mu_{\nu \rho} Y^\rho
  \label{eq:covariant-derivative-vector}
\end{equation}
Note that $ \tensor{Y}{^\mu_{;\nu}} $ is the $ \mu^{\text{th}} $ component of $ \na_\nu Y $, while $ \tensor{Y}{^\mu_{,\nu}} $ is the partial derivative of $ Y^{\mu} $ along $ \pa_{\nu} $.

The covariant derivative coincides with other derivatives on $ \cm $: it can be shown that $ \na_X f = \mathcal{L}_X f = X(f) $ and $ \na_\mu f = \pa_\mu f $. On $ \xm $, however, $ \na_X $ and $ \mathcal{L}_X $ are distinct: while $ \na_X = X^\mu \na_\mu $, there is no way to write the same relation for $ \mathcal{L}_X $, for it depends not only on $ X $ but on its first derivative too. The covariant derivative is thus the natural generalization of the partial derivative to curved manifolds.

\begin{proposition}{}{gamma-non-tens}
  $ \Gamma^\mu_{\rho \nu} $ are not components of a tensor.
\end{proposition}

\begin{proofbox}
  \begin{proof}
    Given the basis transformation $ \tilde{e}_\nu = \tensor{A}{^\mu _\nu} e_\mu $, with $ A $ and invertible matrix (if they are both coordinate basis, then $ \tensor{A}{^\mu _\nu} = \frac{\pa x^\mu}{\pa y^\nu} $), the components of a (1,2) tensor must transform as:
    \begin{equation*}
      \tensor{\tilde{T}}{^\mu _\rho _\nu} = \tensor{(A^{-1})}{^\mu _\tau} \tensor{A}{^\sigma _\rho} \tensor{A}{^\lambda _\nu} \tensor{T}{^\tau _\sigma _\lambda}
    \end{equation*}
    In the new basis:
    \begin{equation*}
      \begin{split}
        \tilde{\Gamma}^\mu_{\rho \nu} \tilde{e}_\mu
        &= \na_{\tilde{e}_\rho}\tilde{e}_\nu = \na_{\tensor{A}{^\sigma _\rho} e_\sigma} (\tensor{A}{^\lambda _\nu} e_{\nu}) = \tensor{A}{^\sigma _\rho} \na_{e_\sigma} (\tensor{A}{^\lambda _\nu} e_\lambda) \\
        &= \tensor{A}{^\sigma _\rho} \tensor{A}{^\lambda _\nu} \Gamma^\tau_{\sigma \lambda} e_\tau + \tensor{A}{^\sigma _\rho} e_\lambda \pa_\sigma \tensor{A}{^\lambda _\nu} = \left[ \tensor{A}{^\sigma _\rho} \tensor{A}{^\lambda _\nu} \Gamma^\tau_{\sigma \lambda} + \tensor{A}{^\sigma _\rho} \pa_\sigma \tensor{A}{^\tau _\nu} \right] e_\tau \\
        &= \left[ \tensor{A}{^\sigma _\rho} \tensor{A}{^\lambda _\nu} \Gamma^\tau_{\sigma \lambda} + \tensor{A}{^\sigma _\rho} \pa_\sigma \tensor{A}{^\tau _\nu} \right] \tensor{(A^{-1})}{^\mu _\tau} \tilde{e}_\mu
      \end{split}
    \end{equation*}
    Thus, there is a second term proportional to $ \pa A $ which deviates from the transformation law:
    \begin{equation*}
      \tilde{\Gamma}^\mu_{\rho \nu} = \tensor{(A^{-1})}{^\mu _\tau} \tensor{A}{^\sigma _\rho} \left[ \tensor{A}{^\lambda _\nu} \Gamma^\tau_{\sigma \lambda} + \pa_\sigma \tensor{A}{^\tau _\nu} \right]
    \end{equation*}
    which means that $ \Gamma^\mu_{\rho \nu} $ is not a tensor.
  \end{proof}
\end{proofbox}

\subsection{Covariant derivative of tensors}

First of all, it is necessary to elucidate how the covariant derivative acts on 1-forms. Given a 1-form $ \omega $, the 1-form $ \na_X \omega $ is defined by its action on vector fields. By Leibniz rule:
\begin{equation*}
  \na_X (\omega(Y)) = (\na_X \omega)(Y) + \omega(\na_X Y)
\end{equation*}
Recalling that $ \omega(Y) $ is a function, $ \na_X (\omega(Y)) = X(\omega(Y)) $, therefore:
\begin{equation}
  (\na_X \omega)(Y) = X(\omega(Y)) - \omega(\na_X Y)
  \label{eq:cov-der-form-aux}
\end{equation}
Expressing it in coordinates:
\begin{equation*}
  \begin{split}
    X^\mu (\na_\mu \omega)_\nu Y^\nu
    &= X^\mu \pa_\mu (\omega_\nu Y^\nu) - \omega_\nu X^\mu \left[ \pa_\mu Y^\nu + \Gamma^\nu_{\mu \rho} Y^\rho \right] \\
    &= X^\mu \left[ \pa_\mu \omega_\rho - \Gamma^\nu_{\mu \rho} \omega_{\nu} \right] Y^\rho
  \end{split}
\end{equation*}
Crucially, the $ \pa Y $ terms cancel out, allowing to define $ \na_X \omega $ without referencing $ Y $:
\begin{equation}
  (\na_\mu \omega)_\rho = \pa_\mu \omega_\rho - \Gamma^\nu_{\mu \rho} \omega_\nu
\end{equation}
Using the same notation as for vector fields $ (\na_\mu \omega)_\rho \equiv \na_\mu \omega_\rho \equiv \omega_{\rho ; \mu} $:
\begin{equation}
  \omega_{\rho ;\mu} = \omega_{\rho,\mu} - \Gamma^\nu_{\mu \rho} \omega_\nu
\end{equation}
This kind of argument can be extended to a general $ (p,q) $ tensor field:
\begin{equation}
  \begin{split}
    \na_\rho \tensor{T}{^{\mu_1} ^\dots ^{\mu_p} _{\nu_1} _\dots _{\mu_q}} = \pa_\rho \tensor{T}{^{\mu_1} ^\dots ^{\mu_p} _{\nu_1} _\dots _{\nu_q}}
    &+ \Gamma^{\mu_1}_{\rho \sigma} \tensor{T}{^\sigma ^{\mu_2} ^\dots ^{\mu_p} _{\nu_1} _\dots _{\nu_q}} + \dots + \Gamma^{\mu_p}_{\rho \sigma} \tensor{T}{^{\mu_1} ^\dots ^{\mu_{p-1}} ^\sigma _{\nu_1} _\dots _{\nu_q}} \\
    &- \Gamma^\sigma_{\rho \nu_1} \tensor{T}{^{\mu_1} ^\dots ^{\mu_p} _\sigma _{\nu_2} _\dots _{\nu_q}} - \dots - \Gamma^\sigma_{\rho \nu_q} \tensor{T}{^{\mu_1} ^\dots ^{\mu_p} _{\nu_1} _\dots _{\nu_{q-1}} _\sigma}
  \end{split}
\end{equation}
The pattern is clear: for each upper index $ \mu $ there is a $ + \Gamma^\mu_{\rho \sigma} T^\sigma $ term, while for each lower index $ \nu $ there is $ - \Gamma^\sigma_{\rho \nu} T_\sigma $ term. Furthermore, it is necessary to generalize the comma-notation: for example, $ \tensor{X}{^\mu _; _\nu _\rho} \equiv \na_\rho \na_\nu X^\mu $, so the rightmost index is the one acting first.

\subsubsection{Torsion and curvature}

Even though the connection is not a tensor, it is used to construct two important tensors.

\begin{definition}{Torsion tensor}{}
  The \bcdef{torsion} is a (1,2) tensor defined on $ \lm{1} \times \xm \times \xm $ as:
  \begin{equation}
    \tors(\omega,X,Y) \defeq \omega (\na_X Y - \na_Y X - [X,Y])
  \end{equation}
  Alternatively, the torsion can be viewed as a map $ \tors : \xm \times \xm \rightarrow \xm $ such that:
  \begin{equation}
    \tors(X,Y) = \na_X Y - \na_Y X - [X,Y]
  \end{equation}
\end{definition}

\begin{definition}{Curvature tensor}{}
  The \bcdef{curvature} is a (1,3) tensor defined on $ \lm{1} \times \xm \times \xm \times \xm $ as:
  \begin{equation}
    \riem(\omega,Z,X,Y) \defeq \omega (\na_X \na_Y Z - \na_Y \na_X Z - \na_{[X,Y]} Z)
    \label{eq:riemann-def}
  \end{equation}
  Alternatively, the curvature can be viewed as a map from $ \xm \times \xm $ to the space of differential operators on $ \xm $ such that:
  \begin{equation}
    \riem(X,Y) = \na_X \na_Y - \na_Y \na_X - \na_{[X,Y]}
    \label{eq:riemann-operator-def}
  \end{equation}
\end{definition}

The fact that these are indeed tensors, i.e. they are linear in each argument, can be shown by direct calculation, recalling that $ [fX,Y] = f[X,Y] - Y(f) X $.

\begin{proposition}{}{}
  On the coordinate basis $ \{\pa_\mu\} $ and $ \{\dd x^\mu\} $ the torsion components are:
  \begin{equation}
    \tensor{T}{^\rho _\mu _\nu} = \Gamma^\rho_{\mu \nu} - \Gamma^\rho_{\nu \mu}
    \label{eq:torsion-comp}
  \end{equation}
\end{proposition}

\begin{proofbox}
  \begin{proof}
    By direct calculation:
    \begin{equation*}
      \begin{split}
        \tensor{T}{^\rho _\mu _\nu}
        &= T(\dd x^\rho, \pa_\mu, \pa_\nu) = \dd x^\mu (\na_\mu \pa_\nu - \na_\nu \pa_\mu - [\pa_\mu, \pa_\nu]) \\
        &= \dd x^\mu (\pa_\mu \pa_\nu - \Gamma^\sigma_{\mu \nu} \pa_\sigma - \pa_\nu \pa_\mu + \Gamma^\sigma_{\nu \mu} \pa_\sigma) \\
        &= \left[ \Gamma^\sigma_{\mu \nu} - \Gamma^\sigma_{\nu \mu} \right] \delta^\mu_\sigma = \Gamma^\rho_{\mu \nu} - \Gamma^\rho_{\nu \mu}
      \end{split}
    \end{equation*}
    which is the thesis.
  \end{proof}
\end{proofbox}

Interestingly, even though $ \Gamma^\rho_{\mu \nu} $ is not a tensor, its anti-symmetric part $ \Gamma^\rho_{[\mu \nu]} = \frac{1}{2} \tensor{T}{^\sigma_\mu _\nu} $ is. Clearly, the torsion tensor is anti-symmetric in its lower indices, thus for connections which are symmetric in their lower indices the torsion is null: such connections are said to be \bctxt{torsion-free}.

\begin{proposition}{}{}
  On the coordinate basis $ \{\pa_\mu\} $ and $ \{\dd x^\mu\} $ the curvature components are:
  \begin{equation}
    \tensor{R}{^\sigma _\rho _\mu _\nu} = \pa_\mu \Gamma^\sigma_{\nu \rho} - \pa_\nu \Gamma^\sigma_{\mu \rho} + \Gamma^\lambda_{\nu \rho} \Gamma^\sigma_{\mu \lambda} - \Gamma^\lambda_{\mu \rho} \Gamma^\sigma_{\nu \lambda}
    \label{eq:riemann-comp}
  \end{equation}
\end{proposition}

\begin{proofbox}
  \begin{proof}
    By direct calculation:
    \begin{equation*}
      \begin{split}
        \riem(\dd x^\sigma, \pa_\rho, \pa_\mu, \pa_\nu)
        &= \dd x^\sigma (\na_\mu \na_\nu \pa_\rho - \na_\nu \na_\mu \pa_\rho - \na_{[\pa_\mu, \pa_\nu]} \pa_\rho) \\
        &= \dd x^\sigma (\na_\mu \na_\nu \pa_\rho - \na_\nu \na_\mu \pa_\rho) \\
        &= \dd x^\sigma (\na_\mu (\Gamma^\lambda_{\nu \rho} \pa_\lambda) - \na_\nu (\Gamma^\lambda_{\mu \rho} \pa_\lambda)) \\
        &= \dd x^\sigma ((\pa_\mu \Gamma^\lambda_{\nu \rho}) \pa_\lambda + \Gamma^\lambda_{\nu \rho} \Gamma^\tau_{\mu \lambda} \pa_\tau - (\pa_\nu \Gamma^\lambda_{\mu \rho}) \pa_\lambda - \Gamma^\lambda_{\mu \rho} \Gamma^\tau_{\nu \lambda} \pa_{\tau}) \\
        &= \pa_\mu \Gamma^\sigma_{\nu \rho} - \pa_\nu \Gamma^\sigma_{\mu \rho} + \Gamma^\lambda_{\nu \rho} \Gamma^\sigma_{\mu \lambda} - \Gamma^\lambda_{\mu \rho} \Gamma^\sigma_{\nu \lambda}
      \end{split}
    \end{equation*}
    which completes the proof.
  \end{proof}
\end{proofbox}

Clearly, the curvature tensor is anti-symmetric in its last two lower indices, i.e. $ \tensor{R}{^\sigma _\rho _\mu _\nu} = \tensor{R}{^\sigma _\rho _[ _\mu _\nu _]} $, and it can in fact be written as:
\begin{equation}
  \tensor{R}{^\sigma _\rho _\mu _\nu} = 2 \pa_{[\mu} \Gamma^\sigma_{\nu] \rho} + 2 \Gamma^\sigma_{[\mu | \lambda |} \Gamma^\lambda_{\nu]\rho}
\end{equation}

\begin{theorem}[before upper = {\tcbtitle}]{Ricci identity}{}
  \begin{equation}
    2 \na_{[\mu} \na_{\nu]} Z^\sigma = \tensor{R}{^\sigma _\rho _\mu _\nu} Z^\rho - \tensor{T}{^\rho_{\mu \nu}} \na_\rho Z^\sigma
    \label{eq:ricci-id}
  \end{equation}
\end{theorem}

\begin{proofbox}
  \begin{proof}
    By direct calculation:
    \begin{equation*}
      \begin{split}
        \na_{[\mu} \na_{\nu]} Z^\sigma
        &= \pa_{[\mu} (\na_{\nu]} Z^\sigma) + \Gamma^\sigma_{[\mu | \lambda |} \na_{\nu]} Z^\lambda - \Gamma^\rho_{[\mu \nu]} \na_\rho Z^\sigma \\
        &= \pa_{[\mu} \pa_{\nu]} Z^\sigma + (\pa_{[\mu} \Gamma^\sigma_{\nu]\rho}) Z^\rho + (\pa_{[\mu} Z^\rho) \Gamma^\sigma_{\nu]\rho} + \Gamma^\sigma_{[\mu | \lambda |} \pa_{\nu]} Z\lambda + \Gamma^\sigma_{[\mu | \lambda |} \Gamma^\lambda_{\nu] \rho} Z^\rho - \frac{1}{2} \tensor{T}{^\rho_{\mu \nu}} \na_\rho Z^\sigma \\
        &= \left( \pa_{[\mu} \Gamma^\sigma_{\nu]\rho} + \Gamma^\sigma_{[\mu | \lambda |} \Gamma^\lambda_{\nu]\rho} \right) Z^\rho - \frac{1}{2} \tensor{T}{^\rho_{\mu \nu}} \na_\rho Z^\sigma \\
        &= \frac{1}{2} \tensor{R}{^\sigma_{\rho \mu \nu}} Z^\rho - \frac{1}{2} \tensor{T}{^\sigma_{\mu \nu}} \na_\rho Z^\sigma
      \end{split}
    \end{equation*}
    which is the thesis.
  \end{proof}
\end{proofbox}

\subsubsection{Levi--Civita connection}

The discussion on the connection has so far been independent of the metric. Starting to consider it, an important result is the fundamental theorem of Riemannian geometry.

\begin{theorem}{Riemann's theorem}{}
  On a metric manifold $ (\mathcal{M},g) $, there exists a unique torsion-free connection that is compatible with the metric, i.e. for all $ X \in \xm $:
  \begin{equation}
    \na_X g = 0
  \end{equation}
  This is called the \bcth{Levi--Civita connection}.
\end{theorem}

\begin{proofbox}
  \begin{proof}
    WTS uniqueness: suppose such a connection exists. Then, by Leibniz rule:
    \begin{equation*}
      X(g(Y,Z)) = \na_X (g(Y,Z)) = (\na_X g)(Y,Z) + g(\na_X Y, Z) + g(Y, \na_X, Z)
    \end{equation*}
    Since $ \na_X g = 0 $, by cyclic permutations of $ X $, $ Y $ and $ Z $:
    \begin{equation*}
      \begin{split}
        X(g(Y,Z)) &= g(\na_X Y, Z) + g(Y, \na_X Z) \\
        Y(g(Z,X)) &= g(\na_Y Z, X) + g(Z, \na_Y X) \\
        Z(g(X,Y)) &= g(\na_Z X, Y) + g(X, \na_Z Y)
      \end{split}
    \end{equation*}
    Since the connection is torsion-free, $ \na_X Y - \na_Y X = [X,Y] $, thus these equations become:
    \begin{equation*}
      \begin{split}
        X(g(Y,Z)) &= g(\na_Y X, Z) + g(\na_X Z, Y) + g([X,Y], Z) \\
        Y(g(Z,X)) &= g(\na_Z Y, X) + g(\na_Y X, Z) + g([Y,Z], X) \\
        Z(g(X,Y)) &= g(\na_X Z, Y) + g(\na_Z Y, X) + g([Z,X], y)
      \end{split}
    \end{equation*}
    Adding the first two and subtracting the third:
    \begin{equation*}
      \begin{split}
        g(\na_Y X, Z) = \frac{1}{2} [ & X(g(Y,Z)) + Y(g(Z,X)) + Z(g(X,Y)) \\
                                      & - g([X,Y],Z) - g([Y,Z],X) + g([Z,X],Y) ]
      \end{split}
    \end{equation*}
    The metric is non-degenerate, thus this uniquely specifies the connection. By direct calculation it can be shown that it indeed satisfies all the properties of a connection.
  \end{proof}
\end{proofbox}

It is possible to explicit the components of the Levi--Civita connection in terms of the metric.

\begin{proposition}{Christoffel symbols}{}
  On the coordinate basis $ \{\pa_\mu\} $ and $ \{dx^\mu\} $ the Levi--Civita connection's components, called \bcprop{Christoffel symbols}, are:
  \begin{equation}
    \Gamma^\lambda_{\mu \nu} = \frac{1}{2} g^{\lambda \rho} (\pa_\mu g_{\nu \rho} + \pa_\nu g_{\mu \rho} - \pa_\rho g_{\mu \nu})
  \end{equation}
\end{proposition}

\begin{proofbox}
  \begin{proof}
    Recalling that $ [\pa_\mu, \pa_\nu] = 0 $, $ \Gamma^\lambda_{\mu \nu} g_{\lambda \rho} = g(\na_\mu \pa_\nu, \pa_\rho) = \frac{1}{2} (\pa_\nu g_{\mu \rho} + \pa_\mu g_{\nu \rho} - \pa_\rho g_{\mu \nu}) $.
  \end{proof}
\end{proofbox}

\begin{example}{Euclidean space}{}
  In flat space $ \R^n $, endowed with either Euclidean or Minkowski metric, it is always possible to choose Cartesian coordinates, in which case the Christoffel symbols vanish. Being the Riemann tensor a genuine tensor, it therefore will vanish in all possible coordinate systems on $ \R^n $, even in those with $ \Gamma^\rho_{\mu \nu} \neq 0 $: this expresses the flatness of $ \R^n $.
\end{example}

\subsubsection{Einstein tensor}

In addition to the anti-symmetry in its last two indices, the Riemann tensor presents additional subtle symmetries when using the Levi--Civita connection.

\begin{proposition}{Symmetries of the Riemann tensor}{riemann-symmetries}
  On a metric manifold with a Levi--Civita connection:
  \begin{equation}
    R_{\sigma \rho \mu \nu} = - R_{\sigma \rho \nu \mu} = - R_{\rho \sigma \mu \nu} = R_{\mu \nu \sigma \rho}
    \qquad \qquad
    R_{\sigma [\rho \mu \nu]} = 0
    \label{eq:riemann-symm}
  \end{equation}
\end{proposition}

\begin{proofbox}
  \begin{proof}
    Set normal coordinates centered at a point $ p $: then, $ \Gamma^\mu_{\nu \rho} = 0 $ and $ \pa_\mu g^{\lambda \sigma} = 0 $ at that point. At $ p $, the Riemann tensor can be written as:
    \begin{equation*}
      \begin{split}
        R_{\sigma \rho \mu \nu}
        &= g_{\sigma \lambda} \tensor{R}{^\lambda _{\rho \mu \nu}} = g_{\sigma \lambda} \left[ \pa_\mu \Gamma^\lambda_{\nu \rho} - \pa_\nu \Gamma^\lambda_{\mu \rho} \right] \\
        &= \frac{1}{2} \left[ \pa_\mu \left( \pa_\nu g_{\sigma \rho} + \pa_\rho g_{\nu \sigma} - \pa_\sigma g_{\nu \rho} \right) - \pa_\nu \left( \pa_\mu g_{\sigma \rho} + \pa_\rho g_{\mu \sigma} - \pa_\sigma g_{\mu \rho} \right) \right] \\
        &= \frac{1}{2} \left[ \pa_\mu \pa_\rho g_{\nu \sigma} - \pa_\mu \pa_\sigma g_{\nu \rho} - \pa_\nu \pa_\rho g_{\mu \sigma} + \pa_\nu \pa_\sigma g_{\mu \rho} \right]
      \end{split}
    \end{equation*}
    The symmetries are then manifest, and being these tensor equations they are valid in all coordinate systems.
  \end{proof}
\end{proofbox}

An important computation tool is the Bianchi identity.

\begin{theorem}{Bianchi identity}{}
  On a metric manifold with a Levi--Civita connection:
  \begin{equation}
    \na_{[\lambda} R_{\sigma \rho] \mu \nu} = 0
    \quad \iff \quad
    \tensor{R}{^\sigma_{\rho [\mu \nu ; \lambda]}} = 0
    \label{eq:riemann-bianchi}
  \end{equation}
\end{theorem}

\begin{proofbox}
  \begin{proof}
    The two equations are equivalent, so the proof is of the first one. In normal coordinates $ \na_\mu = \pa_\mu $ at $ p $, so schematically: $ R = \pa \Gamma + \Gamma \Gamma $, thus $ \na R = \pa R = \pa^2 \Gamma + \Gamma \pa \Gamma = \pa^2 \Gamma $. Explicitly:
    \begin{equation*}
      \pa_\lambda R_{\sigma \rho \mu \nu} = \frac{1}{2} \pa_\lambda \left[ \pa_\mu \pa_\rho g_{\nu \sigma} - \pa_\mu \pa_\sigma g_{\nu \rho} - \pa_\nu \pa_\rho g_{\mu \sigma} + \pa_\nu \pa_\sigma g_{\mu \rho} \right]
    \end{equation*}
    Anti-symmetrizing the appropriate indices yields the result.
  \end{proof}
\end{proofbox}

Note that \eeref{eq:riemann-symm}{eq:riemann-bianchi} do not require that the connection is a Levi--Civita connection, but are valid for general torsion-free connections.

\begin{definition}{Ricci tensor and scalar}{}
  On a metric manifold, the \bcdef{Ricci tensor} is defined as:
  \begin{equation}
    R_{\mu \nu} \defeq \tensor{R}{^\rho _{\mu \rho \nu}}
  \end{equation}
  The \bcdef{Ricci scalar} is defined as:
  \begin{equation}
    R \defeq g^{\mu \nu} R_{\mu \nu}
  \end{equation}
\end{definition}

From \eref{eq:riemann-symm}, the Ricci tensor is clearly symmetric: $ R_{\mu \nu} = R_{\nu \mu} $.

\begin{proposition}{}{}
  On a metric manifold:
  \begin{equation}
    \na^\mu R_{\mu \nu} = \frac{1}{2} \na_\nu R
    \label{eq:ricci}
  \end{equation}
\end{proposition}

\begin{proofbox}
  \begin{proof}
    Writing explicitly Bianchi identity:
    \begin{equation*}
      \na_\lambda R_{\sigma \rho \mu \nu} + \na_\sigma R_{\rho \lambda \mu \nu} + \na_\rho R_{\lambda \sigma \mu \nu} = 0
    \end{equation*}
    Contracting with $ g^{\mu \lambda} g^{\rho \nu} $:
    \begin{equation*}
      \na^\mu R_{\mu \sigma} - \na_\sigma R + \na^\nu R_{\nu \sigma} = 0
    \end{equation*}
    which yields the thesis.
  \end{proof}
\end{proofbox}

An important tensor obtained as a combination of the Ricci tensor and scalar is the \bctxt{Einstein tensor}, which encodes information on the curvature of the metric manifold on which it is defined:
\begin{equation}
  G_{\mu \nu} \defeq R_{\mu \nu} - \frac{1}{2} R g_{\mu \nu}
  \label{eq:einstein-def}
\end{equation}
Most importantly, the Einstein tensor is covariantly conserved, as of \eref{eq:ricci}:
\begin{equation}
  \na^\mu G_{\mu \nu} = 0
\end{equation}

\subsubsection{Gauss' theorem}

The divergence theorem (or Gauss' theorem) states that the integral of a total derivative is a boundary term. It is possible to express this theorem on curved manifolds in a convenient way.

\begin{lemma}[before upper = {\tcbtitle}]{}{gauss-lemma}
  \begin{equation}
    \Gamma^\mu_{\mu \nu} = \frac{1}{\sqg} \pa_\nu \sqg
  \end{equation}
\end{lemma}

\begin{proofbox}
  \begin{proof}
    A useful identity for invertible matrices: $ \tr \log \mt{A} = \log \det \mt{A} $. Thus (WLOG $ \det g > 0 $):
    \begin{equation*}
      \begin{split}
        \Gamma^\mu_{\mu \nu}
        &= \frac{1}{2} g^{\mu \rho} \pa_\nu g_{\mu \rho} = \frac{1}{2} \tr (g^{-1} \pa_\nu g) = \frac{1}{2} \tr (\pa_\nu \log g) = \frac{1}{2} \pa_\nu \log \det g = \frac{1}{\sqrt{\det g}} \pa_\nu \sqrt{\det g}
      \end{split}
    \end{equation*}
    which is the thesis.
  \end{proof}
\end{proofbox}

\begin{theorem}{Gauss' theorem}{}
  Given a Riemannian manifold $ (\mathcal{M},g) $, consider a region $ M \subseteq \mathcal{M} $ with boundary $ \pa M $ and let $ n^\mu $ be an outward-pointing unit vector orthogonal to $ \pa M $. Then, for any vector field $ X^\mu $ on $ M $:
  \begin{equation}
    \int_M \dd^n x \sqg\, \na_\mu X^\mu = \int_{\pa M} \dd^{n-1} x \sqrt{\gamma}\, n_\mu X^\mu
    \label{eq:gauss-theorem}
  \end{equation}
  where $ \gamma_{ij} $ is the pull-back of the metric to $ \pa M $ and $ \gamma \equiv \det \gamma_{ij} $.
\end{theorem}

\begin{proofbox}
  \begin{proof}
    From \lref{lemma:gauss-lemma}:
    \begin{equation*}
      \sqg\, \na_\mu X^\mu = \sqg \left( \pa_\mu X^\mu + \Gamma^\mu_{\mu \nu} X^\nu \right) = \sqg \left( \pa_\mu X^\mu + X^\nu \frac{1}{\sqg} \pa_\nu \sqg \right) = \pa_\mu ( \sqg X^\mu )
    \end{equation*}
    The integral becomes:
    \begin{equation*}
      \int_M \dd^n x \sqg\, \na_\mu X^\mu = \int_M \dd^n x\, \pa_\mu (\sqg X^\mu)
    \end{equation*}
    This is the integral of an ordinary partial derivative, so the ordinary divergence theorem applies. To evaluate the integral on the boundary, it is convenient to pick coordinates so that $ \pa M $ is a surface at constant $ x^n $. Moreover, to simplify the proof, the possible metrics will be restricted to $ g_{\mu \nu} = \diag \left( \gamma_{ij}, N^2 \right) $. By usual integration rules:
    \begin{equation*}
      \int_M \dd^n x\, \pa_\mu ( \sqg X^\mu) = \int_{\pa M} \dd^{n-1} x \sqrt{\gamma N^2} X^n
    \end{equation*}
    The unit normal vector is $ n^\mu = (0,\dots,0,\frac{1}{N}) $, so that $ g_{\mu \nu} n^\mu n^\nu = 1 $, therefore $ n_\mu = g_{\mu \nu} n^\nu = (0,\dots,0,N) $. The proof is then concluded because:
    \begin{equation*}
      \int_{\pa M} \dd^{n-1} x\sqrt{\gamma N^2} X^n = \int_{\pa M} \dd^{n-1} x \sqrt{\gamma} n_\mu X^\mu
    \end{equation*}
    which is the thesis.
  \end{proof}
\end{proofbox}

Note that this theorem holds on Lorentzian manifolds too, with the condition that $ \pa M $ must be purely timelike or purely spacelike, ensuring that $ \gamma \neq 0 $ at any point.

\newpage
\section{Parallel transport}

The connection connects tangent spaces, or more generally any tensor vector space, at different points of the manifold: this map is called parallel transport and it's necessary for the definition of differentiation.

\begin{definition}{Parallel transport}{}
  Consider a vector field $ X $ and some associated integral curve $ \gamma $, with coordinates $ x^\mu(\tau) $ such that:
  \begin{equation*}
    X^\mu \big\vert_\gamma = \frac{\dd x^\mu(\tau)}{\dd \tau}
  \end{equation*}
  A tensor field $ T $ is said to be \bcdef{parallely transported} along $ \gamma $ if:
  \begin{equation}
    \na_X T = 0
    \label{eq:parallel-transport}
  \end{equation}
\end{definition}

Suppose that $ \gamma $ connects two points $ p,q \in \mathcal{M} $: \eref{eq:parallel-transport} provides a map from the tensor vector space defined at $ p $ to that defined at $ q $. To illustrate this, consider the parallel transport of a vector field $ Y $:
\begin{equation*}
  X^\nu (\pa_\nu Y^\mu + \Gamma^\mu_{\nu \rho} Y^\rho) = 0
\end{equation*}
Evaluating this equation on $ \gamma $, considering $ Y^\mu = Y^\mu(x(\tau)) $:
\begin{equation}
  \frac{\dd Y^\mu}{\dd \tau} + X^\nu \Gamma^\mu_{\nu \rho} Y^\rho = 0
\end{equation}
These are a set of coupled ODEs, thus, given an initial condition (e.g. at $ \tau = 0 $, i.e. at $ p $), these equations can be solved to find a unique vector at each point along the curve.

Note that the parallel transport depends both on the path (characterized by the vector field $ X $) and on the connection.

\begin{definition}{Geodesic}{}
  Given a vector field $ X $, a \bcdef{geodesic} is an integral curve to $ X $ such that:
  \begin{equation}
    \na_X X = 0
    \label{eq:geodesic-def}
  \end{equation}
\end{definition}

\begin{proposition}{Geodesic equation}{}
  A geodesic is described by:
  \begin{equation}
    \frac{\dd^2 x^\mu}{\dd \tau^2} + \Gamma^\mu_{\nu \rho} \frac{\dd x^\nu}{\dd \tau} \frac{\dd x^\rho}{\dd \tau} = 0
    \label{eq:geodesic-equation}
  \end{equation}
\end{proposition}

\begin{proofbox}
  \begin{proof}
    From the above calculations, along $ \gamma $:
    \begin{equation*}
      0 = \frac{\dd X^\mu}{\dd \tau} + \Gamma^\mu_{\nu \rho} X^\nu X^\rho = \frac{\dd^2 x^\mu}{\dd \tau^2} + \Gamma^\mu_{\nu \rho} \frac{\dd x^\nu}{\dd \tau} \frac{\dd x^\rho}{\dd \tau}
    \end{equation*}
    which is the geodesic equation.
  \end{proof}
\end{proofbox}

For the Levi--Civita connection $ \na_X g = 0 $. If $ Y \in \xm $ is parallely transported along a geodesic associated to $ X \in \xm $, then $ \na_X Y = \na_X X = 0 $, therefore $ \frac{\dd}{\dd \tau} g(X,Y) = 0 $: this ensures that the two tangent vectors always determine the same angles along the geodesic.

\subsection{Normal coordinates}

\begin{theorem}{Equivalence principle}{equivalence}
  Given a Riemannian manifold $ (\mathcal{M},g) $ and $ p \in \mathcal{M} $, in a neighbourhood of $ p $ it's always possible to find coordinates, called \bcth{normal coordinates}, such that $ g_{\mu \nu}(p) = \delta_{\mu \nu} $ and $ g_{\mu \nu , \rho}(p) = 0 $.
\end{theorem}

\begin{proofbox}
  \begin{proof}
    By brute force, consider initial coordinates $ y^\mu $ and find a change of coordinates $ x^\mu $ which satisfy the requirements:
    \begin{equation*}
      \frac{\pa y^\rho}{\pa x^\mu} \frac{\pa y^\sigma}{\pa x^\nu} \tilde{g}_{\rho \sigma} = g_{\mu \nu}
    \end{equation*}
    WLOG $ p $ is the origin of both sets of coordinates, so:
    \begin{equation*}
      y^\rho = \frac{\pa y^\rho}{\pa x^\mu}\bigg\vert_{x = 0} x^\mu + \frac{1}{2} \frac{\pa^2 y^\rho}{\pa x^\mu \pa x^\nu}\bigg\vert_{x = 0} x^\mu x^\nu + \dots
    \end{equation*}
    This, together with the Taylor expansion of $ \tilde{g}_{\rho \sigma} $, can be inserted in the transformation equation of the metric, thus finding a set of PDEs for each power of $ x $, which can be solved to characterize the normal coordinates. For example, the first condition is:
    \begin{equation*}
      \frac{\pa y^\rho}{\pa x^\mu}\bigg\vert_{x = 0} \frac{\pa y^\sigma}{\pa x^\nu}\bigg\vert_{x = 0} \tilde{g}_{\rho \sigma}(p) = \delta_{\mu \nu}
    \end{equation*}
    Given any $ \tilde{g}_{\rho \sigma}(p) $, it's always possible to find $ \frac{\pa y}{\pa x} $ that satisfies this condition. In fact, if $ \dim_{\R} \mathcal{M} = n $, the Jacobian of the transformation has $ n^2 $ independent elements and the equation above puts $ \frac{1}{2} n (n + 1) $ constraints: the remaining free parameters are $ \frac{1}{2} n (n - 1) $, which is precisely the dimension of $ \SOn{n} $, the symmetry group of the flat metric.
    A similar counting shows that $ g_{\mu \nu , \rho}(p) = 0 $ puts $ \frac{1}{2} n^2 (n + 1) $ constraints, precisely the number of independent elements of the Hessian of the transformation.
  \end{proof}
\end{proofbox}

This theorem holds for Lorentzian manifolds too, but the flat metric is now $ \eta_{\mu \nu} $ and its symmetry group is $ \SOn{1, n-1} $. The condition $ g_{\mu \nu , \rho}(p) = 0 $ implies that $ \Gamma^\mu_{\nu \rho}(p) = 0 $, but generally the Christoffel symbols won't vanish away from $ p $. Note, however, that it's not generally possible to ensure the vanishing of second derivatives too: indeed, $ g_{\mu \nu , \rho \sigma}(p) = 0 $ would put $ \frac{1}{4} n^2 (n + 1)^2 $ constraints, but the independent $ \frac{\pa^3 y}{\pa^3 x} $ terms are $ \frac{1}{6} n^2 (n+1) (n+2) $, thus leaving $ \frac{1}{12} n^2 (n^2 - 1) $ free terms: this is precisely the number of independent components of the Riemann tensor, therefore in general it is not possible to pick coordinates as to make the Riemann tensor vanish too.

\subsubsection{Exponential map}

\begin{figure}
  \centering
  \includegraphics[width = 0.70 \textwidth]{exp-map.png}
  \caption{Visualization of the exponential map.}
  \label{fig:exp-map}
\end{figure}

A simple way to construct normal coordinates is the following: given a tangent vector $ X_p \in T_p \mathcal{M} $, there is a unique affinely parametrized geodesic through $ p $ with tangent vector $ X_p $ at $ p $; then, any point $ q $ in the neighbourhood of $ p $ is labelled by the coordinates of the geodesic that takes from $ p $ to $ q $ a fixed amount of time.

Analytically, introducing a coordinate system (not necessarily normal) $ \tilde{x}^\mu $ in the neighbourhood of $ p $, an affinely parametrized geodesic solves \eref{eq:geodesic-equation}, with initial conditions $ \frac{\pa \tilde{x}^\mu}{\pa \tau}\vert_{\tau = 0} = \tilde{X}_p^\mu $ and $ \tilde{x}^\mu (\tau = 0) = 0 $ that make the solution unique. The uniqueness of the solution allows to define a map $ \Exp : T_p \mathcal{M} \rightarrow \mathcal{M} $, called the \bctxt{exponential map}, which acts as follows: given $ X_p \in T_p \mathcal{M} $, construct the appropriate geodesic as above and follow it for a fixed affine distance, conventionally $ \tau = 1 $, to get a new point $ q \in \mathcal{M} $. See Fig. \figref{fig:exp-map} for visual aid. Obviously, there may be points which cannot be reached from $ p $ by geodesics, or there may be tangent vectors $ X_p $ for which $ \Exp $ is ill-defined: in General Relativity, this occurs when spacetime has singularities, but these are not relevant issues.

Pick a basis $ \{e_\mu\} $ of $ T_p \mathcal{M} $. Then $ \Exp : T_p \mathcal{M} \ni X^\mu e_\mu \mapsto q \in \mathcal{M} $, thus it is possible to assign coordinates in the neighbourhood of $ p $ such that $ x^\mu(q) = X^\mu $: these are the normal coordinates. To show this, note that if $ \{e_\mu\} $ is orthonormal, then the geodesics will point in orthogonal directions, ensuring that $ g_{\mu \nu}(p) = \delta_{\mu \nu} $. Now, fix a point $ q $ associated to a given tangent vector $ X_p \in T_p \mathcal{M} $: this means that $ q $ is at distance $ \tau = 1 $ from $ p $ along the given geodesic. Note that the geodesic equation is homogeneous in $ \tau $, thus in general $ \Exp : \tau X_p \mapsto x^\mu(\tau) = \tau X^\mu $, which means that geodesics take a simple form in these coordinates:
\begin{equation*}
  x^\mu (\tau) = \tau X^\mu
\end{equation*}
Being these geodesics, they must solve \eref{eq:geodesic-equation}, that is:
\begin{equation*}
  \Gamma^\mu_{\nu \rho}(x(\tau)) X^\nu X^\rho = 0
\end{equation*}
which holds at any point along the geodesic, i.e. at any $ \tau \in \R^+ $. At most points $ x(\tau) $, this equation only holds for those choices of $ X^\mu $ tangent to the geodesics. However, at $ x(0) = 0 $, i.e. at $ p $, it must hold for any tangent vector: this means that $ \Gamma^\mu_{(\nu \rho)}(p) = 0 $ which, for a torsion-free connection, ensures that $ \Gamma^\mu_{\nu \rho}(p) = 0 $. But vanishing Christoffel symbols imply a vanishing first derivative of the metric: for the Levi--Civita connection $ 2g_{\mu \sigma} \Gamma^\sigma_{\nu \rho} = g_{\mu \nu, \rho} + g_{\mu \rho, \nu} - g_{\nu \rho, \mu} $, thus symmetrizing $ (\mu \nu) $ cancels the last two terms, leaving an identity that, evaluated at $ p $, gives $ g_{\mu \nu, \rho}(p) = 0 $. Hence, these are indeed normal coordinates.

\subsubsection{Equivalence principle}

Normal coordinates are conceptually important in General Relativity: an observer at point $ p $ who parametrizes their immediate surroundings using coordinates constructed by geodesics will experience a locally flat metric. This is precisely Einstein's equivalence principle: any free-falling observer, performing local experiments, will not experience a gravitational field. The formal definition of free-falling observer is an observer which follows geodesics, while the local lack of gravitational field means $ g_{\mu \mu}(p) = \eta_{\mu \nu} $. In this context, normal coordinates are called \bctxt{local inertial frame}.

To understand what ``local" means, note that there is a way to distinguish whether a gravitational field is present at $ p $: a non-vanishing Riemann tensor. This depends on the second derivatives of the metric, which in general will be non-vanishing. However, to measure the effects of the Riemann tensor, one needs to compare the results of experiments at $ p $ and at a nearby point $ q $: this is a non-local observation.

\subsection{Curvature and torsion}

\begin{figure}[!b]
  \centering
  \includegraphics[width = 0.90 \textwidth]{parall-tr.png}
  \caption{Parallel transport along different paths.}
  \label{fig:par-tr}
\end{figure}

With reference to \figref{fig:par-tr}, consider a tangent vector $ Z_p \in T_p \mathcal{M} $ and two vector fields $ X,Y \in \xm : [X,Y] = 0 $, i.e. they are linearly independent. Construct two curves $ \gamma, \gamma' $ as in figure, both leading to a point $ r \in \mathcal{M} $ which, for simplicity, is close to $ p $. It is possible to impose normal coordinates centered at $ p $ such that $ x^\mu = (\tau, \sigma, \dots) $, so that $ X = \frac{\pa}{\pa \tau} $ and $ Y = \frac{\pa}{\pa \sigma} $: then $ x^\mu(p) = (0,0,0, \dots) $, $ X^\mu(q) = (\delta \tau, 0, 0, \dots) $, $ x^\mu(s) = (0, \delta \sigma, 0, \dots) $ and $ x^\mu(r) = (\delta \tau, \delta \sigma, 0, \dots)$, with $ \delta \tau, \delta \sigma $ small.\\
First, parallel transport $ Z_p $ along $ X $ to $ Z_q $, so that $ Z^\mu $ solves:
\begin{equation*}
  \frac{\dd Z^\mu}{\dd \tau} + X^\nu \Gamma^\mu_{\nu \rho} Z^\rho = 0
\end{equation*}
In normal coordinates $ \Gamma^\mu_{\nu \rho}(p) = 0 $, thus $ \frac{\dd Z^\mu}{\dd \tau}\big\vert_{\tau = 0} = 0 $ and the Taylor expansion is:
\begin{equation*}
  \begin{split}
    Z_q^\mu
    &= Z_p^\mu + \frac{\delta \tau^2}{2} \frac{\dd^2 Z^\mu}{\dd \tau^2}\bigg\vert_{\tau = 0} + \smo(\delta\tau^3) \\
    &= Z_p^\mu - \frac{\delta\tau^2}{s} \left[ X^\nu Z^\rho \frac{\dd \Gamma^\mu_{\nu \rho}}{\dd \tau} + \frac{\dd X^\nu}{\dd \tau} Z^\rho \Gamma^\mu_{\nu \rho} + X^\nu \frac{\dd Z^\rho}{\dd \tau} \Gamma^\mu_{\nu \rho} \right]_{\tau = 0} + \smo(\delta\tau^3) \\
    &= Z_p^\mu - \frac{\delta\tau^2}{2} X^\nu Z^\rho \frac{\dd \Gamma^\mu_{\nu \rho}}{\dd \tau}\bigg\vert_{\tau = 0} + \smo(\delta\tau^3) = Z_p^\mu - \frac{\delta\tau^2}{2} \left[ X^\nu X^\sigma Z^\rho \Gamma^\mu_{\nu \rho , \sigma} \right]_p + \smo(\delta\tau^3)
  \end{split}
\end{equation*}
where $ \frac{\dd}{\dd \tau} = X^\sigma \pa_\sigma $. Now, $ Z_q $ needs to be parallely transported along $ Y $ to $ Z_r $, but this time $ \frac{\dd Z^\mu}{\dd \sigma}\big\vert_{\sigma = 0} $ doesn't vanish, in general. From the parallel transport equation:
\begin{equation*}
  \frac{\dd Z^\mu}{\dd \sigma}\bigg\vert_{\sigma = 0} = - \left[ Y^\nu Z^\rho \Gamma^\mu_{\nu \rho} \right]_q = - \left[ Y^\nu Z^\rho X^\sigma \Gamma^\mu_{\nu \rho , \sigma} \right]_p \delta \tau + \smo(\delta\tau^2)
\end{equation*}
The expansions of $ Y^\nu $ and $ Z^\rho $ at leading order multiply $ \Gamma^\mu_{\nu \rho}(p) = 0 $, thus only contribute to higher order terms. Next order in $ \delta \sigma $:
\begin{equation*}
  \frac{\dd^2 Z^\mu}{\dd \sigma^2}\bigg\vert_{\sigma = 0} = - \left[ \left( \frac{\dd Y^\nu}{\dd \sigma} Z^\rho + Y^\nu \frac{\dd Z^\rho}{\dd \sigma} \right) \Gamma^\mu_{\nu \rho} + Y^\nu Z^\rho \frac{\dd \Gamma^\mu_{\nu \rho}}{\dd \sigma} \right]_q = - \left[ Y^\nu Y^\sigma Z^\rho \Gamma^\mu_{\nu \rho , \sigma} \right]_p + \smo(\delta\tau)
\end{equation*}
The complete expansion thus is:
\begin{equation*}
  \begin{split}
    Z_r^\mu
    &= Z_q^\mu - \left[ Y^\nu Z^\rho X^\sigma \Gamma^\mu_{\nu \rho , \sigma} \right]_p \delta\tau \delta\sigma - \frac{1}{2} \left[ Y^\nu Y^\sigma Z^\rho \Gamma^\mu_{\nu \rho , \sigma} \right]_p \delta\sigma^2 + \smo(\delta^3) \\
    &= Z_p^\mu - \frac{1}{2} \Gamma^\mu_{\nu \rho , \sigma}(p) \left[ X^\nu X^\sigma Z^\rho \delta\tau^2 + 2 Y^\nu Z^\rho X^\sigma \delta\tau \delta\sigma + Y^\nu Y^\sigma Z^\rho \delta\sigma^2 \right] + \smo(\delta^3)
  \end{split}
\end{equation*}
Parallel transport along $ \gamma' $ leads to a similar expression (exchange of $ X $ and $ Y $):
\begin{equation*}
  {Z'_r}^\mu = Z_p^\mu - \frac{1}{2} \Gamma^\mu_{\nu \rho , \sigma}(p) \left[ Y^\nu Y^\sigma Z^\rho \delta\sigma^2 + 2 X^\nu Z^\rho Y^\sigma \delta\sigma \delta\tau + X^\nu X^\sigma Z^\rho \delta\tau^2 \right] + \smo(\delta^3)
\end{equation*}
The difference between the parallely transported tangent vectors to leading order is:
\begin{equation*}
  \Delta Z_r^\mu = Z_r^\mu - {Z'_r}^\mu = - \left[ \Gamma^\mu_{\nu \rho , \sigma} - \Gamma^\mu_{\sigma \rho , \nu} \right]_p \left[ Y^\nu Z^\rho X^\sigma \right]_p \delta\tau \delta\sigma + \smo(\delta^3)
\end{equation*}
Recalling that $ \Gamma^\mu_{\nu \rho}(p) = 0 $, it is possible to write:
\begin{equation}
  \Delta Z_r^\mu = - \left[ \tensor{R}{^\mu_{\rho \sigma \nu}} Y^\nu Z^\rho X^\sigma \right]_p \delta\sigma \delta\tau + \smo(\delta^3)
  \label{eq:deviation}
\end{equation}
It would be possible to evaluate the expression at $ r $ too, as it would differ only by higher order terms. Although the calculation was carried in a particular choice of coordinates, \eref{eq:deviation} is a tensor relation, therefore it must hold in all coordinate systems: the Riemann tensor thus determines the path dependence of parallel transport.

\subsubsection{Torsion}

\begin{figure}
  \centering
  \includegraphics[width = 0.50 \textwidth]{torsion.png}
  \caption{Visualization of torsion.}
  \label{fig:torsion}
\end{figure}

Consider two tangent vectors $ X_p,Y_p \in T_p \mathcal{M} $ and a coordinate system $ x^\mu $ such that $ X_p = X^\mu \pa_\mu $ and $ Y_p = Y^\mu \pa_\mu $. If $ p : x^\mu $, as in \figref{fig:torsion} construct $ r,s \in \mathcal{M} $ such that $ r : x^\mu + \varepsilon X^\mu $ and $ s : x^\mu + \varepsilon Y^\mu $, with $ \varepsilon $ an infinitesimal parameter. Now, parallel transport $ X_p \in T_p \mathcal{M} $ along the direction of $ Y_p $ to $ X'_s \in T_s \mathcal{M} $ and $ Y_p \in T_p \mathcal{M} $ along $ X_p $ to $ Y'_r \in T_r \mathcal{M} $; their components will be:
\begin{equation*}
  X'_s = \left( X^\mu - \varepsilon \Gamma^\mu_{\nu \rho} Y^\nu X^\rho \right)
  \qquad \qquad
  Y'_r = \left( Y^\mu - \varepsilon \Gamma^\mu_{\nu \rho} X^\nu Y^\rho \right)
\end{equation*}
Repeating this process, starting from point $ s $ and moving along the direction of $ X'_s $, a new point $ q \in \mathcal{M} $ is determined, with coordinates:
\begin{equation*}
  q : x^\mu + \varepsilon \left( X^\mu + Y^\mu \right) - \varepsilon^2 \Gamma^\mu_{\nu \rho} Y^\nu X^\rho
\end{equation*}
Analogously, starting at point $ r $ and moving along the direction of $ Y'_r $, a new point $ t \in \mathcal{M} $ is determined, with coordinates:
\begin{equation*}
  t : x^\mu + \varepsilon \left( X^\mu + Y^\mu \right) - \varepsilon^2 \Gamma^\mu_{\nu \rho} X^\nu Y^\rho
\end{equation*}
If the connection is torsion-free, then $ q \equiv t $. On the other hand, if $ T^\mu_{\nu \rho} \neq 0 $, the parallelogram fails to close, as in \figref{fig:torsion}.

\subsection{Geodesic deviation}
\label{ssec:geo-dev}

\begin{definition}{Tangent and deviation vector fields}{}
  Given a one-parameter family of geodesics $ \{x^\mu(\tau;s)\}_{s \in \R} $ on a manifold $ \mathcal{M} $, the \bcdef{tangent vector field} and the \bcdef{deviation vector field} are defined as:
  \begin{equation}
    X^\mu \defeq \frac{\pa x^\mu}{\pa \tau}\bigg\vert_s
    \qquad \qquad
    S^\mu \defeq \frac{\pa x^\mu}{\pa s}\bigg\vert_\tau
  \end{equation}
\end{definition}
The meaning of these vector fields is evident: the tangent vector field fixes a particular geodesics (i.e. a particular $ s $) and assigns at each point of the geodesic its tangent vector, while the deviation vector field fixes a particular value of the affine parameter $ \tau $ and assigns at each point with this value a vector which takes to a nearby geodesic (at the same $ \tau $).

\begin{figure}
  \centering
  \includegraphics[width = 0.70 \textwidth]{geo-dev.png}
  \caption{A one-parameter family of geodesics generated by $ X $.}
  \label{fig:geo-dev}
\end{figure}

The family of geodesics sweeps a surface embedded in the manifold, so there is freedom in the choice of coordinates $ s $ and $ \tau $. In particular, it's always possible to pick them so that $ X = \frac{\pa}{\pa \tau} $ and $ S = \frac{\pa}{\pa s} $, in order for them to be linearly independent: $ [X,S] = 0 $, as in \figref{fig:geo-dev}. Consider a torsion-free connection, so that:
\begin{equation*}
  \na_X S - \na_S X = [X,S] = 0
  \quad \implies \quad
  \na_X \na_X S = \na_X \na_S X = \na_S \na_X X + \riem(X,S) X
\end{equation*}
But $ X $ is tangent to geodesics, so from \eref{eq:geodesic-def} $ \na_X X = 0 $ and:
\begin{equation}
  \na_X \na_X S = \riem(X,S) X
\end{equation}
Restricting to an integral curve $ \gamma $ of the vector field $ X $, i.e. $ X^\mu \vert_\gamma = \frac{\dd x^\mu}{\dd \tau} $, the covariant derivative along $ \gamma $ becomes:
\begin{equation*}
  \na_X\big\vert_\gamma = X^\mu\big\vert_\gamma \na_\mu = \frac{\dd x^\mu}{\dd \tau} \na_\mu \equiv \frac{D}{D\tau}
\end{equation*}
Hence, in index notation, the change of the deviation vector along the geodesic is expressed as:
\begin{equation}
  \frac{D^2 S^\mu}{D\tau^2} = \tensor{R}{^\mu _{\nu \rho \sigma}} X^\rho S^\sigma X^\nu
  \label{eq:geodesic-deviation}
\end{equation}
This can be interpreted as the relative acceleration of neighbouring geodesics, and it is determined by the Riemann tensor. Experimentally, such geodesic deviations are called \bctxt{tidal forces}, which are the non-local manifestation of a gravitational field as a curvature of spacetime.
